% Team Note of PS akgwi
% These codes should be guaranteed, fast enough, short and easy to type.

% twocolumn
% \documentclass[landscape, 8pt, a4paper, oneside, twocolumn]{extarticle}

\documentclass[landscape, 8pt, a4paper, oneside]{extarticle}
\usepackage{kotex}
\usepackage{amssymb}
\usepackage{amsmath}
\usepackage{import}
\usepackage{enumitem}
\usepackage{multicol}
\usepackage{titlesec}

\usepackage{teamnote}

\teamnote{Dankook University}{ThreeMinuteSearch}{gokim06, auaahks, wnstjd9288}{ICPC 2026 Seoul regional}

% smaller section title
% \titleformat{\subsection}{\normalfont\normalsize\bfseries}{\thesubsection}{1em}{}

\titlespacing*{\section}{0pt}{0pt}{0pt}
\titlespacing*{\subsection}{0pt}{0pt}{0pt}

\ShowUsage
\ShowComplexity
\HideAuthor

\begin{document}

\begin{multicols*}{2}

\maketitlepage

% break column
% \vfill\null\columnbreak

% Make Pagebreak if you want.
% \pagebreak

\section{정렬 / 이분 탐색 관련}

\Algorithm{정렬 비교 함수 구현}{Point 구조체의 연산자 오버로딩 또는 Compare 함수로 정렬 기준 지정}{-}{cpp}{code/sorting/custom_sort_comparator.cpp}

\Algorithm{좌표 압축}{원소의 대소관계를 유지하면서 [0, N) 범위의 수로 압축}{O(N \log N)}{cpp}{code/sorting/coordinate_compression.cpp}


\section{그래프/트리 (Easy)}

\Algorithm{위상 정렬}{AddEdge로 간선 추가 후 TopSort 호출}{O(V + E)}{cpp}{code/graph/topological_sort.cpp}

\Algorithm{최단 거리 - Floyd Warshall}{Init으로 초기화, AddEdge로 간선 추가 후 Run 호출}{O(V^3)}{cpp}{code/graph/floyd_warshall.cpp}

\Algorithm{최단 거리 - Dijkstra}{음수 가중치가 없는 그래프에서 시작점 s와 도착점 t를 인자로 Dijkstra 호출}{O(E \log E)}{cpp}{code/graph/dijkstra.cpp}

\Algorithm{최단 거리 - Bellman Ford}{AddEdge로 간선 추가 후 시작점 s를 인자로 Run 호출}{O(VE)}{cpp}{code/graph/bellman_ford.cpp}

\Algorithm{유니온 파인드 + 최소 신장 트리(Kruskal)}{Find와 Merge로 유니온 파인드 수행, 간선을 정렬한 뒤 Kruskal 수행}{UF: 연산마다 O(\log N), MST: O(E \log E)}{cpp}{code/graph/disjoint_set_union_kruskal.cpp}


\section{자료구조}

\Algorithm{세그먼트 트리}{Set으로 원소를 변경하고 Sum으로 [l, r] 구간의 합 계산}{O(\log N)}{cpp}{code/data_structure/segment_tree.cpp}

\Algorithm{세그먼트 트리 + 레이지 프로퍼게이션}{RangeAdd로 구간에 값을 더하고 RangeSum으로 구간 합 계산}{O(\log N)}{cpp}{code/data_structure/lazy_segment_tree.cpp}

\Algorithm{Fenwick Tree}{range_update로 구간에 값을 더하고 rsq로 구간 합 계산}{O(\log N)}{cpp}{code/data_structure/fenwick_tree.cpp}

\Algorithm{Policy Based Data Structures - Ordered Set}{find_by_order로 k번째 원소 탐색, order_of_key로 특정 원소보다 작은 원소 개수 계산}{O(\log N)}{cpp}{code/data_structure/policy_based_ordered_set.cpp}

\Algorithm{퍼시스턴트 세그먼트 트리}{사용 전에 init(root[0], s, e) 호출}{O(\log N)}{cpp}{code/data_structure/persistent_segment_tree.cpp}


\section{그래프/트리 (Hard)}

\Algorithm{SCC - Kosaraju}{정방향 DFS로 종료 순서를 저장한 뒤 역방향 DFS로 SCC 구성}{O(V + E)}{cpp}{code/graph/strongly_connected_components_kosaraju.cpp}

\Algorithm{BCC - Tarjan}{간선 추가 후 다른 작업을 수행하기 전에 tarjan 호출}{O(V + E)}{cpp}{code/graph/biconnected_components_tarjan.cpp}

\Algorithm{최대 유량 - Dinic}{maximum_flow(s, t)로 s에서 t로의 최대 유량 계산}{O(V^2 E), 모든 간선의 용량이 1이면 O(\min(V^{2/3}, E^{1/2})E)}{cpp}{code/graph/dinic_maximum_flow.cpp}

\Algorithm{MCMF}{-}{-}{cpp}{code/graph/minimum_cost_maximum_flow.cpp}

\Algorithm{이분 매칭 - Hopcroft Karp}{-}{O(E\sqrt{V})}{cpp}{code/graph/hopcroft_karp_matching.cpp}

\Algorithm{Heavy Light Decomposition}{정점 가중치 또는 간선 가중치에 대한 전처리 후 경로 쿼리 수행}{전처리 O(N), 쿼리 O(T(N)\log N)}{cpp}{code/tree/heavy_light_decomposition.cpp}

\Algorithm{Centroid Decomposition}{-}{-}{cpp}{code/tree/centroid_decomposition.cpp}


\section{수학}

\Algorithm{선형 시간 체, 곱셈적 함수 전처리}{Sieve로 소수 및 곱셈적 함수 전처리}{전처리 O(P), 쿼리 O(\log P)}{cpp}{code/math/linear_sieve_multiplicative_functions.cpp}

\Algorithm{확장 유클리드 알고리즘}{확장 유클리드 알고리즘으로 최대공약수와 계수 계산}{O(\log \max(a,b))}{cpp}{code/math/extended_euclidean_algorithm.cpp}

\Algorithm{중국인의 나머지 정리}{확장 유클리드 알고리즘을 이용해 연립 합동식 계산}{O(k \log m)}{cpp}{code/math/chinese_remainder_theorem.cpp}

\Algorithm{이항 계수를 소수로 나눈 나머지}{-}{-}{cpp}{code/math/lucas_theorem_binomial_mod_prime.cpp}

\Algorithm{빠른 소수 판별, 소인수분해 - Miller Rabin, Pollard Rho}{처음에 Sieve() 호출}{IsPrime: O(\log^2 N), Factorize: 약 O(N^{1/4})}{cpp}{code/math/miller_rabin_pollard_rho.cpp}

\Algorithm{가우스 소거법 - RREF, 랭크, 행렬식, 역행렬}{-}{O(N^3)}{cpp}{code/math/gaussian_elimination_rref.cpp}

\Algorithm{다항식 곱셈(FFT)}{multiply로 두 다항식의 계수 계산, multiply\_mod로 모듈러 연산 수행}{O(N \log N)}{cpp}{code/math/fast_fourier_transform.cpp}


\section{문자열}

\Algorithm{문자열 해싱}{build로 문자열 전처리, get으로 부분 문자열의 해시값 계산}{build: O(N), get: O(1)}{cpp}{code/string/string_hashing.cpp}

\Algorithm{문자열 매칭 - KMP}{KMP(s, p)로 문자열 s에서 패턴 p가 등장하는 위치 반환}{GetFail: O(|P|), KMP: O(|S| + |P|)}{cpp}{code/string/knuth_morris_pratt.cpp}

\Algorithm{가장 긴 팰린드롬 부분 문자열 - Manacher}{Manacher로 각 문자를 중심으로 하는 최장 팰린드롬의 반경 반환}{O(N)}{cpp}{code/string/manacher_palindrome.cpp}

\Algorithm{문자열 매칭 - Z}{Z[i]로 S[0:N]과 S[i:N]의 최장 공통 접두사 길이 계산}{O(N)}{cpp}{code/string/z_algorithm.cpp}

\Algorithm{접미사 배열}{SuffixArray로 접미사 배열과 LCP 배열 계산}{O(N \log N)}{cpp}{code/string/suffix_array.cpp}


\section{계산 기하}

\Algorithm{2차원 계산 기하 템플릿 + CCW}{-}{-}{cpp}{code/geometry/geometry_template_ccw.cpp}

\Algorithm{360도 각도 정렬}{-}{-}{cpp}{code/geometry/polar_angle_sorting.cpp}

\Algorithm{다각형 넓이}{-}{-}{cpp}{code/geometry/polygon_area.cpp}

\Algorithm{선분 교차 판정}{-}{-}{cpp}{code/geometry/segment_intersection.cpp}

\Algorithm{다각형 내부 판별}{PointInPolygon으로 점이 다각형 내부 또는 경계 위에 있는지 판별}{O(N)}{cpp}{code/geometry/point_in_polygon.cpp}

\Algorithm{볼록 껍질 - Graham Scan}{ConvexHull로 모든 점을 포함하는 가장 작은 볼록 다각형 계산}{O(N \log N)}{cpp}{code/geometry/convex_hull_graham_scan.cpp}

\Algorithm{가장 먼 두 점 - Rotating Calipers}{반시계 방향으로 정렬된 볼록 다각형을 입력받아 가장 먼 두 점 계산}{O(N)}{cpp}{code/geometry/rotating_calipers_diameter.cpp}

\Algorithm{볼록 다각형 내부 판별}{반시계 방향으로 정렬된 볼록 다각형에서 점이 내부 또는 경계 위에 있는지 판별}{O(\log N)}{cpp}{code/geometry/point_in_convex_polygon.cpp}


\section{기타}

\Algorithm{가장 긴 증가하는 부분 수열(LIS)}{LIS 함수로 가장 긴 증가하는 부분 수열 반환}{O(N \log N)}{cpp}{code/etc/longest_increasing_subsequence.cpp}

\Algorithm{삼분 탐색}{Check 함수의 최솟값을 갖는 위치 탐색}{-}{cpp}{code/etc/ternary_search.cpp}

\Algorithm{C++ 랜덤, GCC 확장, 비트마스킹 트릭}{-}{-}{cpp}{code/etc/random_gcc_extensions_bit_tricks.cpp}

\Algorithm{빠른 입력(Fast Input from STDIN)}{STDIN으로부터 빠르게 입력}{-}{cpp}{code/etc/fast_input.cpp}


\section{Other Stuffs}

\Algorithm{Sum Over Subsets Dynamic Programming}{-}{-}{cpp}{code/etc/sum_over_subsets_dp.cpp}

\Algorithm{Edmond's Algorithm (Directed-MST)}{-}{O(VE)}{cpp}{code/graph/directed_minimum_spanning_tree_edmonds.cpp}

\Algorithm{LCA - ETT + RMQ}{전처리 후 두 정점의 LCA 쿼리 수행}{preprocessing: O(N \log N)}{cpp}{code/tree/lca_euler_tour_rmq.cpp}

\Algorithm{LCA - Binary Lifting}{-}{-}{cpp}{code/tree/lca_binary_lifting.cpp}

\Algorithm{Eulerian Path/Circuit}{-}{-}{cpp}{code/graph/eulerian_path_circuit.cpp}

\Algorithm{Maximum Subarray Segment Tree}{-}{-}{cpp}{code/data_structure/maximum_subarray_segment_tree.cpp}

\Algorithm{Offline Queries - Mo's Algorithm}{-}{O(Q \log Q + (N + Q)N)}{cpp}{code/etc/mo_algorithm.cpp}

\Algorithm{2-Dimensional Segment Tree - Recursive}{-}{-}{cpp}{code/data_structure/two_dimensional_segment_tree_recursive.cpp}

\Algorithm{2-Dimensional Segment Tree - Iterative}{-}{-}{cpp}{code/data_structure/two_dimensional_segment_tree_iterative.cpp}

\Algorithm{2-Dimensional Fenwick Tree}{-}{-}{cpp}{code/data_structure/two_dimensional_fenwick_tree.cpp}

% \section{Notes}

% \subsection{Triangles/Trigonometry} % from Ryute Teamnote
% \input{code/Notes/Triangle}

% \subsection{Calculus, Newton's Method}
% \input{code/Notes/Calculus}

% \subsection{Zeta/Mobius Transform}
% \input{code/Notes/ZetaMobius}

% \subsection{Generating Function}
% \input{code/Notes/GeneratingFunction}

% \subsection{Counting}
% \input{code/Notes/Counting}

% \subsection{Faulhaber's Formula ($\sum_{k=1}^{n} k^c$)}
% \input{code/Notes/Faulhaber}

% \subsection{About Graph Degree Sequence}
% \input{code/Notes/DegreeSequence}

% \subsection{Burnside, Grundy, Pick, Hall, Simpson, Area of Quadrangle, Fermat Point, Euler, Pythagorean}
% \input{code/Notes/Theorems}

% \subsection{About Graph Minimum Cut}
% \input{code/Notes/MinCut}

% \subsection{Matrix with Graph(Kirchhoff, Tutte, LGV)}
% \input{code/Notes/GraphMatrix}

% \subsection{About Graph Matching(Graph with $\vert V \vert \leq 500$)}
% \input{code/Notes/Matching}

% \subsection{Checklist}
% \input{code/Notes/Checklist}

\end{multicols*}

\end{document}
